\clearpage
\section*{September 10, 2026: Resolve documented unit counts and project boundaries}
\textit{Agent-drafted review and implementation of supported corrections.}

The unit review identifies 137 omitted apartments at 1920 Turnbull Avenue. HDC's June 2026 authorization lists 228 apartments: 142 studios, 54 one-bedroom and 32 two-bedroom apartments. The developer describes 137 supportive apartments, 90 affordable apartments and one superintendent apartment. The saved 91-unit measure therefore covers only part of the proposed residential program. Using 228 raises the linked parent's total from 298 to 435 units. At Boyland, the December 2025 zoning sheet confirms 28 apartments at 670 and 86 at 688, retaining 114 despite a conflicting 29-unit DOB field.

Astoria Cove's December 2023 geotechnical report identifies four buildings as Phase 1 on the four lots associated with our January 13, 2022 filings. These establish one parent containing 240, 138, 108 and 90 units, previously divided between two parents. The merged total is 576 in the saved Housing Database vintage. Later proposals have different counts; neither this merger nor the Turnbull correction reconstructs the designs at their original filing dates. Documented proposed-design counts now take priority where explicitly adjudicated, while both original source counts remain available.

The boundary review retains Linden and Rockaway as common developments spanning separately financed phases. Contemporary reporting corroborates Nostrand's three buildings. The Coney Island neighbors remain separate, as do the Ninth Avenue candidates and Wilson/Boston under the documented baseline decisions. Shared ownership or an award covering several buildings alone does not establish a common development, and these economic groupings do not establish common legal wage assessment.

Seven groups remain unresolved: Wharf, 769 East 232nd, Grand Concourse, Myrtle, Beach 30th, Mt Hope and the possible across-street Boone extension. Their conflicting unit measures or incomplete boundaries are recorded in section 12 of the detailed case review. Available records did not supply the missing filing schedules or decisive project evidence. Estimation remains on hold.

The rebuilt panel has 2,733 parents, down from 2,734, and the complete-characteristic A/B review has 66 multi-constituent parents, down from 67. All 8,054 membership rows remain. Only Turnbull's selected count changes; raw unit fields, filing and cohort dates, and refiling indicators are unchanged. Every field of unrelated parents matches the preceding build. The neighborhood screen falls from 1,210 to 1,206 pairs after the Astoria merger.

\textit{Reproduction:} The committed decisions in \path{tasks/parent_opportunities_manual/output} feed the parent constructor and data-panel Make targets. The casebook and source URLs are in \path{tasks/audits/audit_estimation_parent_links}; public PDFs are acquired by its source-download task. Data vintages remain HDB 25Q4 and DOB July 2026, supplemented by dated documentary evidence. Branch \path{simplify_empirical_tasks}, changes after 44c2d2a. Data and audit builds were checked without running estimation.
